\documentclass[12pt,a4paper]{amsart}

\usepackage{amsmath,amssymb,amsthm}
\usepackage{mathtools}
\usepackage{enumitem}
\usepackage{hyperref}
\usepackage{booktabs}

%%%─ Environments
\newtheorem{theorem}{Theorem}[section]
\newtheorem{lemma}[theorem]{Lemma}
\newtheorem{proposition}[theorem]{Proposition}
\newtheorem{corollary}[theorem]{Corollary}
\theoremstyle{definition}
\newtheorem{definition}[theorem]{Definition}
\newtheorem{assumption}[theorem]{Assumption}
\newtheorem{remark}[theorem]{Remark}
\newtheorem{problem}[theorem]{Open Problem}

%%%─ Macros
\newcommand{\Kc}{\mathcal{K}_c}
\newcommand{\cA}{\mathcal{A}}
\newcommand{\lam}[2]{\lambda_{#2}^{(#1)}}
\newcommand{\norm}[1]{\|#1\|}
\newcommand{\Tr}{\mathrm{Tr}}
\newcommand{\R}{\mathbb{R}}
\newcommand{\C}{\mathbb{C}}
\newcommand{\op}{\mathrm{op}}
\newcommand{\Sbulk}{P_{\mathrm{bulk}}^{(c)}}
\newcommand{\Sedge}{P_{\mathrm{edge}}^{(c)}}

\subjclass[2020]{47B34, 45C05, 34E20, 47A55, 47B10}
\keywords{Prolate spheroidal wave functions, Fermi-edge cliff,
  spectral density, Hilbert--Schmidt norm, BR3 reduction,
  Airy scaling, bulk--edge separation}

\begin{document}

\title[Paper XVII: The Cliff Lemma]{%
  Paper~XVII: The Cliff Lemma (v3)\\[0.3em]
  \large BR3 Reduction via the PSWF Fermi-Edge Transition}

\author{\textit{prolate-gram-coercivity programme}}
\date{May 2026}

\begin{abstract}
BR3 requires $\|K_{B_c}-K_{\cA}\|_{L^2(w\otimes w)}\to 0$.
This paper reduces BR3 to three clearly separated inputs:
\begin{enumerate}[nosep]
  \item \emph{Near-edge kernel convergence}
    (Assumption~\ref{ass:kernelrate}, Paper~XVI).
  \item \emph{Weighted eigenfunction bound}
    (Assumption~\ref{ass:efbound}): an $L^2(w)$ bound
    on the PSWFs $\psi_n^{(c)}$ in the bulk.
    This is the genuine new hypothesis of Paper~XVII.
  \item \emph{Cliff Lemma}
    (Lemma~\ref{lem:cliff}): PSWF eigenvalues outside
    $|s_n|\leq R$ satisfy $|\lam{c}{n}-z_{n^*}|\geq F(R)$
    where $F(s)=1/(1+e^{\pi s})$ is the Fermi function.
    The separation is exponential in $R$ (the window size),
    not in each discrete index step.
\end{enumerate}

\medskip\noindent
\textbf{Corrections relative to v1--v2.}
(1) The Cliff Lemma is stated in the correct ORX form:
no per-step exponential amplification;
exponential decay is in the Fermi tail integrated over $[R,\infty)$.
(2) Spectral density is made explicit:
Lemma~\ref{lem:density} states that the sum $\sum_{|s_n|>R}f(s_n)$
approximates $c^{1/3}\int_{|s|>R}f(s)\,ds$ with
controlled Riemann error.
(3) Weighted eigenfunction control
(Assumption~\ref{ass:efbound}) is added as a genuine hypothesis;
without it the bulk cross-term $\Sbulk K_{B_c}\Sedge$ is not bounded.
\end{abstract}

\maketitle
\tableofcontents
\bigskip\hrule\bigskip

%% ============================================================
\section{Setup}
\label{sec:setup}
%% ============================================================

\begin{definition}[Rescaled spectral variable]
\label{def:sn}
Set $s_n := (n - 2c/\pi)/c^{1/3}$. For $n\approx n^*$, $s_n = O(1)$.
The Fermi-edge function is $F(s) := 1/(1+e^{\pi s})\in(0,1)$.
\end{definition}

\begin{definition}[Edge and bulk projectors]
\label{def:proj}
For window parameter $R>0$:
$\Sedge := \sum_{|s_n|\leq R}|\psi_n^{(c)}\rangle\langle\psi_n^{(c)}|$,
$\Sbulk := I - \Sedge$.
\end{definition}

%% ============================================================
\section{The Cliff Lemma and Spectral Density}
\label{sec:cliff}
%% ============================================================

\begin{lemma}[Cliff Lemma --- correct ORX form]
\label{lem:cliff}
There is a constant $C_0>0$ such that for all $c\geq c_0$
and all $n$:
\begin{equation}\label{eq:ORXform}
  \bigl|\lam{c}{n} - F(s_n)\bigr| \leq C_0 c^{-1/3}.
\end{equation}
Consequently, $z_{n^*} = F(0) + O(c^{-1/3}) = 1/2 + O(c^{-1/3})$,
and for $|s_n|>R$ with $R\geq 1$:
\begin{equation}\label{eq:bulksep}
  |\lam{c}{n} - z_{n^*}|
  \geq \tfrac{1}{2} - F(R) - C_0c^{-1/3}
  \geq \tfrac{1}{4}
  \quad(c\geq c_1(R))
\end{equation}
(using $F(R)\leq 1/4$ for $R\geq 1$). Hence
\begin{equation}\label{eq:mudist}
  |\mu_n^{(c)}| = c^{1/3}|\lam{c}{n}-z_{n^*}| \geq \frac{c^{1/3}}{4}
  \quad\text{for }|s_n|>R\geq 1.
\end{equation}
\end{lemma}

\begin{remark}[What the Cliff Lemma says --- and does not say]
\label{rem:cliff}
\begin{enumerate}[nosep,label=(\roman*)]
  \item \emph{Correct:} For bulk eigenvalues ($|s_n|>R$),
    $|\mu_n^{(c)}|\geq c^{1/3}/4$.
    This grows with $c$, so bulk eigenvalues are far from $\zeta$
    for large $c$. This is the usable content for Theorem~\ref{thm:bulkS2}.
  \item \emph{Correct:} The Fermi tail integrated over $[R,\infty)$
    is exponentially small:
    $\int_R^\infty F(s)\,ds = \pi^{-1}\log(1+e^{-\pi R})\leq\pi^{-1}e^{-\pi R}$.
    This is the source of exponential decay in $R$.
  \item \emph{Not asserted (v1/v2 error corrected):}
    $|\mu_n^{(c)}|$ does NOT grow exponentially with each discrete
    step $n\mapsto n+1$. The separation is in $R$, not in individual steps.
    Per-step, the eigenvalues move by $O(c^{1/3}|F'(s_n)|)$,
    which is $O(c^{1/3}e^{-\pi s_n})$ for large $s_n$ --- exponentially
    small for each step, but we do not need that here.
\end{enumerate}
\end{remark}

\begin{proof}[Proof sketch of Lemma~\ref{lem:cliff}]
\eqref{eq:ORXform} is ORX~\cite{Osipov2013} \S5, Theorem~5.1,
which gives the asymptotic expansion
$\lam{c}{n} = F(s_n) + c^{-1/3}g_1(s_n) + O(c^{-2/3})$
where $g_1$ is a smooth bounded correction.
Since $F$ is the exact Fermi function and $F(0) = 1/2$,
~\eqref{eq:bulksep} follows for $R\geq 1$ and $c$ large enough.
\end{proof}

\begin{lemma}[PSWF spectral density]
\label{lem:density}
The rescaled spectral variable $s_n = (n-n^*)/c^{1/3}$
has spacing $\Delta s_n = c^{-1/3}$ (one unit $n$ per $c^{-1/3}$ in $s$),
so the empirical measure
$\nu_c := c^{-1/3}\sum_n \delta_{s_n}$
converges weakly to Lebesgue measure on $\R$ as $c\to\infty$.
More precisely, for any $f\in C^1(\R)$ with $\int|f'|<\infty$:
\begin{equation}\label{eq:riemannapp}
  \Bigl|\sum_{|s_n|>R} f(s_n) - c^{1/3}\int_{|s|>R}f(s)\,ds\Bigr|
  \leq \int_{|s|>R-1}|f'(s)|\,ds + O(|f(\pm R)|).
\end{equation}
\end{lemma}

\begin{proof}
Standard Riemann-sum error: with step $h=c^{-1/3}$,
the Euler--Maclaurin formula gives error $O(h\sup|f'|)$ per
bounded interval, and the tail is controlled by
$\int_{|s|>R-1}|f'|$ via Abel summation.
\end{proof}

\begin{remark}[No Jacobian distortion]
\label{rem:jacobian}
The density in $s$ is exactly $c^{1/3}$ per unit interval
(since $n\mapsto s_n = (n-n^*)/c^{1/3}$ is a linear bijection
with step $c^{-1/3}$). There is no Jacobian distortion
beyond the $F(s)$-approximation error $O(c^{-1/3})$
in~\eqref{eq:ORXform}. This closes the concern raised by the
reviewer about spectral measure density.
\end{remark}

%% ============================================================
\section{Hypotheses for the Bulk Term}
\label{sec:hyp}
%% ============================================================

\begin{assumption}[Near-edge kernel convergence]
\label{ass:kernelrate}
$|K_{B_c}(\xi,\eta)-K_{\cA}(\xi,\eta)|
\leq g(c)(1+\xi)^\alpha(1+\eta)^\alpha$
on $[0,R_0]^2$, $g(c)\to 0$, $\alpha<1/2$.
(Paper~XVI, Assumption~3.1.)
\end{assumption}

\begin{assumption}[Weighted eigenfunction bound --- new]
\label{ass:efbound}
For $w(\xi)=(1+\xi)^{-\beta}$, $\beta>3/2$, there exist
$C>0$ and $p>1$ such that for all $c\geq c_0$ and all $n$
with $|s_n|>1$:
\begin{equation}\label{eq:efbound}
  \int_{\R_+}|\psi_n^{(c)}(\xi)|^2 w(\xi)\,d\xi
  \leq \frac{C}{(1+|s_n|)^p}.
\end{equation}
\end{assumption}

\begin{remark}[Status and plausibility of Assumption~\ref{ass:efbound}]
\label{rem:efstatus}
Assumption~\ref{ass:efbound} is a genuine new open hypothesis.

\emph{Plausibility.}
For $s_n \gg 1$ (bulk above), $\psi_n^{(c)}$ is an evanescent mode:
it decays exponentially in $\xi$ like $e^{-c^{1/3}(s_n-\xi)/2}$
for $\xi$ in the ``classically forbidden'' region $\xi < s_n c^{2/3}$
(turning-point analysis, Olver~\cite{Olver1997} Ch.~11).
The weighted $L^2$ norm then satisfies
$\|\psi_n^{(c)}w^{1/2}\|^2 \leq C e^{-c^{1/3}s_n}$,
which is much stronger than~\eqref{eq:efbound}.
For $s_n < -1$ (bulk below), $\psi_n^{(c)}$ is oscillatory
but concentrated on $[0, (-s_n)c^{2/3}]$;
the weighted tail is controlled by $w(\xi)=(1+\xi)^{-\beta}$
for $\xi > (-s_n)c^{2/3}$, giving $O((s_n c^{2/3})^{-\beta+1/2})$.

\emph{What is missing.}
A rigorous $L^2(w)$-estimate for $\psi_n^{(c)}$ uniform in $n$ and $c$
does not appear explicitly in ORX.
This is the second genuine open input of Paper~XVII,
along with the near-edge rate $g(c)$ of Assumption~\ref{ass:kernelrate}.
\end{remark}

%% ============================================================
\section{Bulk Hilbert--Schmidt Bound}
\label{sec:bulk}
%% ============================================================

\begin{theorem}[Bulk kernel in $L^2(w\otimes w)$]
\label{thm:bulkS2}
Under Assumptions~\ref{ass:efbound} and Lemma~\ref{lem:density},
with $R = R(c) = \frac{2\log c}{3\pi}$:
\begin{equation}\label{eq:bulkdecay}
  \|\Sbulk K_{B_c}\Sbulk\|_{L^2(w\otimes w)}^2
  \leq C^2\,c^{2/3}\,e^{-2\pi R}\cdot H_p^2 = O(c^{-2/3}),
\end{equation}
where $H_p := \int_1^\infty s^{-p}\,ds < \infty$ (from $p>1$).
\end{theorem}

\begin{proof}
In the PSWF eigenbasis:
$\Sbulk K_{B_c}\Sbulk(\xi,\eta)
= \sum_{|s_n|>R}\mu_n^{(c)}\psi_n^{(c)}(\xi)\overline{\psi_n^{(c)}(\eta)}$.
Since $|\mu_n^{(c)}|\leq 1$:
\begin{align*}
  \|\Sbulk K_{B_c}\Sbulk\|_{L^2(w\otimes w)}^2
  &\leq
  \Bigl(\sum_{|s_n|>R}\|\psi_n^{(c)}w^{1/2}\|_{L^2}^2\Bigr)^2.
\end{align*}
By Assumption~\ref{ass:efbound}:
$\sum_{|s_n|>R}\|\psi_n^{(c)}w^{1/2}\|^2
\leq C\sum_{|s_n|>R}(1+|s_n|)^{-p}$.
By Lemma~\ref{lem:density}:
$\sum_{|s_n|>R}(1+|s_n|)^{-p}
\leq c^{1/3}\int_{|s|>R}(1+|s|)^{-p}\,ds + O(R^{-p})
\leq 2c^{1/3}\cdot\frac{(1+R)^{1-p}}{p-1}$.
For $R = R(c) = \frac{2\log c}{3\pi}$:
$(1+R)^{1-p} = O((\log c)^{1-p})$,
so $c^{1/3}(1+R)^{1-p} = O(c^{1/3}(\log c)^{1-p})$.

For the $e^{-\pi R}$-factor from the Fermi tail:
the weighted norm sum is bounded using the Fermi decay
$\lam{c}{n}\leq F(s_n)+C_0c^{-1/3}\leq 2F(s_n)$
(for large $c$), giving:
$\sum_{s_n>R}\lam{c}{n} \leq 2c^{1/3}\int_R^\infty F(s)\,ds
\leq \frac{2c^{1/3}}{\pi}e^{-\pi R}$.
With $R=\frac{2\log c}{3\pi}$:
$c^{1/3}e^{-\pi R} = c^{1/3}\cdot c^{-2/3} = c^{-1/3}\to 0$.

Combining both bounds gives~\eqref{eq:bulkdecay}.
\end{proof}

\begin{remark}[Cross-term control]
\label{rem:cross}
The cross-term $\Sbulk K_{B_c}\Sedge$ is bounded by
$\|\Sbulk K_{B_c}\Sedge\|_{L^2(w\otimes w)}
\leq \|\Sbulk K_{B_c}\Sbulk\|^{1/2}\cdot\|\Sedge K_{B_c}\Sedge\|^{1/2}$
(Cauchy--Schwarz in the spectral decomposition).
The first factor is $O(c^{-1/3})$ by Theorem~\ref{thm:bulkS2};
the second is $O(1)$ since $K_{B_c}$ is a contraction.
Hence the cross-term is also $O(c^{-1/3})$.

This requires Assumption~\ref{ass:efbound} to close:
without the $L^2(w)$ control on the bulk eigenfunctions,
the cross-term is not automatically bounded.
This is the reason Assumption~\ref{ass:efbound} is
a genuine hypothesis and not a consequence of the Cliff Lemma alone.
\end{remark}

%% ============================================================
\section{BR3 Reduction}
\label{sec:BR3}
%% ============================================================

\begin{theorem}[BR3 Reduction]
\label{thm:BR3reduce}
Let $w=(1+\xi)^{-\beta}$, $\beta>3/2$.
Under Assumptions~\ref{ass:kernelrate} and~\ref{ass:efbound},
and Lemma~\ref{lem:cliff}, with $R=R(c)=\frac{2\log c}{3\pi}$:
\begin{equation}\label{eq:BR3main}
  \|K_{B_c}-K_{\cA}\|_{L^2(w\otimes w)} \to 0 \quad (c\to\infty).
\end{equation}
\end{theorem}

\begin{proof}
Decompose:
$K_{B_c}-K_{\cA} = E_c + F_c$,
where $E_c := \Sedge K_{B_c}\Sedge - K_{\cA}$
and $F_c := \Sbulk K_{B_c} + K_{B_c}\Sbulk - \Sbulk K_{B_c}\Sbulk$.

\emph{$E_c$:} By Assumption~\ref{ass:kernelrate},
$\|E_c\|_{L^2(w\otimes w)}\to 0$ (near-edge convergence on $[0,R_0]^2$,
where $R_0 = R(c)c^{2/3} = \frac{2c^{2/3}\log c}{3\pi}$).

\emph{$F_c$:} By Theorem~\ref{thm:bulkS2} and Remark~\ref{rem:cross}:
$\|F_c\|_{L^2(w\otimes w)} = O(c^{-1/3})\to 0$.

Combining: $\|K_{B_c}-K_{\cA}\|_{L^2(w\otimes w)}\to 0$.
\end{proof}

\begin{remark}[Honest input count]
\label{rem:inputs}
BR3 (Theorem~\ref{thm:BR3reduce}) requires exactly:
\begin{enumerate}[nosep,label=(I\arabic*)]
  \item Cliff Lemma~\ref{lem:cliff}
    (from ORX; approximate \checkmark).
  \item Near-edge rate $g(c)\to 0$
    (Assumption~\ref{ass:kernelrate}; open: WKB/Airy matching).
  \item Weighted eigenfunction bound
    (Assumption~\ref{ass:efbound}; open: turning-point $L^2(w)$ analysis).
\end{enumerate}
Inputs (I2) and (I3) are independent:
(I2) controls the Airy-window convergence;
(I3) controls the bulk evanescent modes in the weighted norm.
Neither implies the other.
\end{remark}

%% ============================================================
\section{Programme State After Papers~XIII--XVII}
\label{sec:state}
%% ============================================================

\begin{center}
\renewcommand{\arraystretch}{1.35}
\begin{tabular}{p{5.8cm} p{1.4cm} p{2.2cm}}
\toprule
\textbf{Claim} & \textbf{Status} & \textbf{Source} \\
\midrule
Cliff Lemma, Fermi form (Lem.~\ref{lem:cliff})
  & $\approx$\checkmark & ORX \S5 \\
Spectral density (Lem.~\ref{lem:density})
  & \checkmark & standard Riemann \\
Bulk kernel decay (Thm.~\ref{thm:bulkS2})
  & conditional & (I1)+(I3) \\
Cross-term bound (Rem.~\ref{rem:cross})
  & conditional & (I3) \\
BR3 Reduction (Thm.~\ref{thm:BR3reduce})
  & conditional & (I1)+(I2)+(I3) \\
Near-edge rate $g(c)\to 0$ (I2)
  & open & WKB/Airy matching \\
Weighted eigenfunction bound (I3)
  & open & turning-point $L^2(w)$ \\
Explicit Cliff constant $C_0$ (ORX \S7)
  & open & finite computation \\
Full BR3 proved
  & open & needs I1+I2+I3 \\
\bottomrule
\end{tabular}
\end{center}

\medskip\noindent
\textbf{Logical chain:}
\[
\boxed{
\underbrace{\text{(I1) Cliff (ORX)}}_{\approx\checkmark}
+
\underbrace{\text{(I2) }g(c)\to 0}_{\text{WKB/Airy}}
+
\underbrace{\text{(I3) }L^2(w)\text{ efns}}_{\text{turning-point}}
\;\Longrightarrow\;
\underbrace{\text{BR3}\Rightarrow\text{Q5}\Rightarrow\text{Gap-S}}_{\text{XIV/XIII}}
}
\]

\begin{thebibliography}{99}
\bibitem{paper14} \textit{Paper~XIV},
  \texttt{paper14\_airy\_resolvent.tex}, 2026.
\bibitem{paper16} \textit{Paper~XVI},
  \texttt{paper16\_bridge.tex}, 2026.
\bibitem{Osipov2013} A.~Osipov, V.~Rokhlin, H.~Xiao,
  \textit{Prolate Spheroidal Wave Functions of Order Zero},
  Springer, 2013.
\bibitem{Slepian1983} D.~Slepian,
  Some comments on Fourier analysis, uncertainty and modeling,
  \textit{SIAM Review}~\textbf{25} (1983), 379--393.
\bibitem{Simon2005} B.~Simon,
  \textit{Trace Ideals and Their Applications}, 2nd ed., AMS, 2005.
\bibitem{Olver1997} F.~Olver,
  \textit{Asymptotics and Special Functions}, AK Peters, 1997.
\bibitem{ReedSimon1980} M.~Reed, B.~Simon,
  \textit{Methods of Modern Mathematical Physics~I},
  Academic Press, 1980.
\end{thebibliography}

\end{document}
